%==============================================================================
%  simple-auto — Precontrato de Servicio (Sesión Privada)
%------------------------------------------------------------------------------
%  Documento PRECONTRACTUAL vinculante. Versión en español (versión que rige).
%  Compilar con pdflatex / xelatex / lualatex.
%
%  CÓMO RELLENARLO
%  ---------------
%  Edita ÚNICAMENTE los \newcommand del bloque "DATOS RELLENABLES" de abajo.
%  El resto del documento toma esos valores automáticamente. Para generarlo por
%  lead desde un script, sustituye los valores entre llaves antes de compilar.
%
%  Los campos del PROVEEDOR están MOCKEADOS — reemplázalos por tus datos reales
%  (nombre fiscal, NIF, domicilio fiscal) cuando los tengas.
%==============================================================================
\documentclass[11pt,a4paper]{article}

\usepackage[utf8]{inputenc}
\usepackage[T1]{fontenc}
\usepackage[spanish,es-noquoting]{babel}
\usepackage[a4paper,margin=2.4cm]{geometry}
\usepackage{lmodern}
\usepackage{enumitem}
\usepackage{titlesec}
\usepackage{xcolor}
\usepackage{fancyhdr}
\usepackage{lastpage}
\usepackage[hidelinks]{hyperref}
\usepackage{array}

\definecolor{accent}{HTML}{D0021B}

%==============================================================================
%  DATOS RELLENABLES  (lo único que debes editar)
%==============================================================================

% --- Proveedor (MOCK — sustituir por tus datos fiscales reales) -------------
\newcommand{\provNombre}{[NOMBRE FISCAL DEL PROVEEDOR]}
\newcommand{\provNif}{[NIF/DNI DEL PROVEEDOR]}
\newcommand{\provDomicilio}{[DOMICILIO FISCAL], Las Palmas de Gran Canaria}
\newcommand{\provEmail}{vmat.icaza@gmail.com}
\newcommand{\provMarca}{simple-auto}

% --- Cliente (se rellena por lead) ------------------------------------------
\newcommand{\cliNombre}{[NOMBRE Y APELLIDOS DEL CLIENTE]}
\newcommand{\cliNif}{[DNI/NIF DEL CLIENTE]}
\newcommand{\cliDomicilio}{[DOMICILIO DEL CLIENTE]}
\newcommand{\cliEmail}{[CORREO DEL CLIENTE]}

% --- Condiciones económicas y fechas ----------------------------------------
\newcommand{\precioBase}{99~\euro}              % precio ordinario
\newcommand{\precioLanz}{49,50~\euro}           % precio de lanzamiento (50\%)
\newcommand{\lanzHasta}{1 de agosto de 2026}    % vigencia del descuento
\newcommand{\fechaContrato}{[FECHA]}            % fecha de aceptación
\newcommand{\fechaSesion}{[FECHA DE LA SESIÓN]} % a concretar tras la reserva

\usepackage{eurosym}

%==============================================================================
%  Estilo
%==============================================================================
\titleformat{\section}{\large\bfseries\color{black}}{\thesection.}{0.6em}{}
\titlespacing*{\section}{0pt}{1.4em}{0.5em}
\setlist[enumerate]{itemsep=2pt,topsep=3pt}
\setlist[itemize]{itemsep=2pt,topsep=3pt}

\pagestyle{fancy}
\fancyhf{}
\renewcommand{\headrulewidth}{0.4pt}
\lhead{\small\textbf{\textcolor{accent}{simple}}\textbf{auto} — Precontrato de servicio}
\rhead{\small\thepage\,/\,\pageref{LastPage}}
\cfoot{\scriptsize\itshape Documento precontractual. Versión en español, que rige a todos los efectos.}

\newcommand{\fieldrule}{\rule{5cm}{0.4pt}}

\begin{document}

%------------------------------------------------------------------------------
\begin{center}
  {\LARGE\bfseries Precontrato de prestación de servicios}\\[2pt]
  {\large Sesión privada de creación y publicación de sitio web}\\[6pt]
  {\small\color{accent}\textbf{simple}\color{black}\textbf{auto}}
\end{center}

\vspace{0.6em}
\noindent En Las Palmas de Gran Canaria, a \fechaContrato.

%------------------------------------------------------------------------------
\section{Partes}

\noindent\textbf{El Proveedor:} \provNombre, con NIF \provNif, domicilio en
\provDomicilio, correo electrónico \href{mailto:\provEmail}{\provEmail}, que opera
bajo la marca «\provMarca».

\smallskip
\noindent\textbf{El Cliente:} \cliNombre, con DNI/NIF \cliNif, domicilio en
\cliDomicilio, correo electrónico \cliEmail.

\smallskip
\noindent Ambas partes se reconocen capacidad legal suficiente para contratar y
acuerdan lo siguiente.

%------------------------------------------------------------------------------
\section{Objeto del contrato}

El Proveedor prestará al Cliente \textbf{una (1) sesión privada guiada}, en remoto
o presencial según se acuerde, cuyo objeto es \textbf{crear y dejar publicado en
internet un sitio web} para el Cliente, y \textbf{enseñar al Cliente a mantenerlo y
mejorarlo por sí mismo} con herramientas gratuitas. La sesión tiene una duración
estimada de hasta dos (2) horas.

%------------------------------------------------------------------------------
\section{Alcance: qué incluye la sesión}

Durante la sesión, el Proveedor realizará junto con el Cliente, como mínimo:
\begin{enumerate}[label=\alph*)]
  \item Publicación de un sitio web funcional en una dirección de internet real,
        en línea y visitable al finalizar la sesión, mediante una plataforma de
        alojamiento gratuita de terceros (Vercel).
  \item Configuración de un repositorio de código (GitHub) y su conexión con la
        plataforma de alojamiento, de modo que cada cambio se publique
        automáticamente.
  \item Ayuda en la adquisición y conexión de un \textbf{nombre de dominio propio}
        del Cliente (p.\,ej. tunombre.com). El coste del dominio (orientativo, unos
        20~\euro/año) lo abona el Cliente directamente al registrador, \textbf{no al
        Proveedor}.
  \item Redacción conjunta de los textos del sitio (titular, secciones y llamada a
        la acción) durante la sesión.
  \item Un formulario de contacto que envía los mensajes al correo del Cliente.
  \item Configuración del sitio para que sea adaptable a móvil, tablet y escritorio
        e indexable por buscadores.
  \item Instalación en el equipo del Cliente de herramientas gratuitas de edición
        asistida por IA (p.\,ej. GitHub Desktop y editores con IA gratuitos), y
        formación práctica para que el Cliente modifique y mejore el sitio por sí
        mismo tras la sesión.
  \item Entrega de una grabación de la sesión.
  \item Acompañamiento durante la sesión hasta dejar el sitio publicado.
\end{enumerate}

%------------------------------------------------------------------------------
\section{Límites del servicio: qué NO incluye}

Para evitar equívocos, el servicio \textbf{no} comprende:
\begin{itemize}
  \item Aplicaciones web complejas, áreas privadas, inicios de sesión, comercio
        electrónico, bases de datos ni automatizaciones a medida. El entregable es
        una página de presentación («landing») sencilla, limpia y rápida.
  \item El coste del nombre de dominio ni de servicios de pago de terceros que el
        Cliente decida contratar.
  \item Garantía de posicionamiento, número de visitas, ventas o resultados
        comerciales. El Proveedor deja el sitio \emph{indexable}, pero no garantiza
        posiciones en buscadores.
  \item Mantenimiento, alojamiento de pago, copias de seguridad, soporte continuo
        ni nuevas funcionalidades posteriores, salvo contratación aparte.
  \item La disponibilidad, precios o condiciones de las herramientas gratuitas de
        terceros (Vercel, GitHub, registradores de dominios, editores con IA), que
        se rigen por sus propios términos y pueden cambiar al margen del Proveedor.
  \item Producción de logotipo, fotografía o redacción de contenidos más allá de lo
        elaborado conjuntamente durante la sesión.
\end{itemize}

%------------------------------------------------------------------------------
\section{Precio y forma de pago}

El precio del servicio es de \precioBase. \textbf{Precio de lanzamiento:} hasta el
\lanzHasta\ se aplica un descuento del 50\,\%, quedando el precio en
\precioLanz. El pago se realiza por adelantado mediante la pasarela indicada por el
Proveedor (PayPal) en el momento de la reserva. La sesión se agenda una vez
confirmado el pago.

%------------------------------------------------------------------------------
\section{Garantía de publicación y reembolso}

El Proveedor garantiza que, al finalizar la sesión, el sitio web quedará
\textbf{publicado y accesible en internet}. \textbf{Si por causa imputable al
Proveedor el sitio no quedara publicado, el Cliente tendrá derecho al reembolso
íntegro} de la cantidad abonada.

\smallskip
La garantía no será de aplicación cuando la imposibilidad de publicar derive de
causas ajenas al Proveedor, en particular: falta de colaboración o de los
materiales necesarios por parte del Cliente, ausencia o incidencias en el equipo o
la conexión del Cliente, falta de pago del dominio o de servicios de terceros, o
caídas o cambios de las plataformas gratuitas de terceros.

%------------------------------------------------------------------------------
\section{Obligaciones del Cliente}

El Cliente se compromete a: aportar los textos, imágenes y materiales con la
antelación solicitada; disponer durante la sesión de un equipo con conexión a
internet y permisos de instalación; crear o facilitar las cuentas gratuitas que se
requieran (correo, GitHub, Vercel, registrador de dominio); y agregar al Proveedor
como colaborador del repositorio si fuera necesario para asistencia. El Cliente es
responsable de la veracidad y legalidad de los contenidos que aporte y publique.

%------------------------------------------------------------------------------
\section{Titularidad}

El sitio web, su código, el dominio y las cuentas de terceros quedan en propiedad y
bajo control \textbf{exclusivo del Cliente}. El Proveedor no retiene titularidad ni
exige suscripción obligatoria alguna para que el Cliente conserve y use su sitio.

%------------------------------------------------------------------------------
\section{Cancelación y reprogramación}

El Cliente podrá reprogramar la sesión con al menos 24~horas de antelación. La
cancelación con menos de 24~horas de antelación, o la incomparecencia, podrá
conllevar la pérdida del importe o su aplicación a una nueva fecha, a criterio
razonable del Proveedor.

%------------------------------------------------------------------------------
\section{Protección de datos}

Los datos personales del Cliente se tratarán por el Proveedor con la única finalidad
de gestionar la relación contractual y la prestación del servicio, conforme al
Reglamento (UE) 2016/679 (RGPD) y a la normativa española aplicable. El Cliente
podrá ejercer sus derechos de acceso, rectificación, supresión y oposición
escribiendo a \href{mailto:\provEmail}{\provEmail}.

%------------------------------------------------------------------------------
\section{Naturaleza precontractual y perfeccionamiento}

El presente documento recoge las condiciones del servicio y tiene carácter
\textbf{vinculante} entre las partes. El contrato se perfecciona con la aceptación
de estas condiciones por el Cliente y la confirmación del pago, momento a partir del
cual ambas partes quedan obligadas en sus términos.

%------------------------------------------------------------------------------
\section{Legislación y fuero}

Este precontrato se rige por la legislación española. Para cualquier controversia,
las partes se someten a los Juzgados y Tribunales de Las Palmas de Gran Canaria,
salvo que por la condición de consumidor del Cliente corresponda otro fuero por
imperativo legal.

%------------------------------------------------------------------------------
\vspace{1.2em}
\section*{Aceptación y firmas}
\addcontentsline{toc}{section}{Aceptación y firmas}

\noindent Las partes aceptan el contenido íntegro de este precontrato.

\vspace{2.2em}
\noindent
\begin{tabular}{p{0.46\textwidth} p{0.46\textwidth}}
\fieldrule & \fieldrule \\[2pt]
\textbf{El Proveedor} & \textbf{El Cliente} \\
\provNombre & \cliNombre \\
NIF: \provNif & DNI/NIF: \cliNif \\
\end{tabular}

\vspace{1.4em}
\noindent Fecha de la sesión acordada: \fechaSesion

\end{document}
